\documentclass[12pt,a4paper]{article}
\usepackage[utf8]{inputenc}
\usepackage[spanish]{babel}
\usepackage{amsmath, amssymb, amsthm}
\usepackage{booktabs}
\usepackage{array}
\usepackage{geometry}
\usepackage{graphicx}
\usepackage{hyperref}
\usepackage{xcolor}
\usepackage{enumitem}
\usepackage{fancyhdr}
\usepackage{titlesec}
\usepackage{tcolorbox}
\usepackage{multirow}
\usepackage{float}

\geometry{top=2.5cm, bottom=2.5cm, left=3cm, right=3cm}

\pagestyle{fancy}
\fancyhf{}
\rhead{Econometría -- Heteroscedasticidad}
\lhead{Tabla 2.8 -- Gujarati}
\rfoot{Página \thepage}

\definecolor{azuloscuro}{RGB}{0,70,127}
\definecolor{grisclaro}{RGB}{240,240,245}
\definecolor{verdeclaro}{RGB}{230,245,230}
\definecolor{rojoclaro}{RGB}{255,230,230}

\titleformat{\section}{\large\bfseries\color{azuloscuro}}{}{0em}{}[\titlerule]
\titleformat{\subsection}{\normalsize\bfseries\color{azuloscuro}}{}{0em}{}

\tcbuselibrary{skins,breakable}
\newtcolorbox{cajaresultado}[1][]{
  colback=grisclaro, colframe=azuloscuro, fonttitle=\bfseries,
  breakable, #1
}
\newtcolorbox{cajaconclucion}[1][]{
  colback=verdeclaro, colframe=green!50!black, fonttitle=\bfseries,
  breakable, #1
}
\newtcolorbox{cajacomando}[1][]{
  colback=black!5, colframe=black!30, fonttitle=\bfseries,
  breakable, #1
}

\begin{document}

%===============================================================
\begin{center}
{\LARGE \textbf{Solución Completa -- Problema 2.15}}\\[0.5em]
{\large Tabla 2.8: Gasto en Comida y Gasto Total (55 familias, India)}\\[0.3em]
{\normalsize \textit{Gujarati -- Econometría Básica, Capítulo 2}}
\end{center}
\hrule height 1.5pt
\vspace{1em}

%===============================================================
\section{Descripción del Problema}

Se cuenta con datos de 55 familias de la India sobre su \textbf{Gasto en Comida} (variable dependiente, $Y$) y su \textbf{Gasto Total} (variable independiente, $X$), ambas medidas en rupias. El objetivo es estudiar la relación entre ambas variables, diagnosticar y corregir el problema de \textbf{heteroscedasticidad}.

El modelo de regresión es:
\[
\text{Gasto\_en\_Comida}_i = \beta_0 + \beta_1 \cdot \text{Gasto\_Total}_i + u_i, \quad i = 1, 2, \ldots, 55
\]

%===============================================================
\section{Inciso A: Regresión del Gasto en Comida sobre el Gasto Total}

\subsection{Comandos en Stata}
\begin{cajacomando}[title=Stata -- Regresión MCO]
\texttt{. rename var2 Gasto\_en\_Comida}\\
\texttt{. rename var3 Gasto\_Total}\\
\texttt{. reg Gasto\_en\_Comida Gasto\_Total}
\end{cajacomando}

\subsection{Resultados de la Regresión (MCO)}

\begin{table}[H]
\centering
\caption{Resultados de Regresión: Gasto\_en\_Comida sobre Gasto\_Total}
\begin{tabular}{lcccccc}
\toprule
\textbf{Variable} & \textbf{Coef.} & \textbf{Error Est.} & \textbf{t} & \textbf{P$>|t|$} & \textbf{IC 95\% Inf.} & \textbf{IC 95\% Sup.}\\
\midrule
Gasto\_Total  & 0.4368 & 0.0783 & 5.58 & 0.000 & 0.2797 & 0.5939\\
\_cons        & 94.2088 & 50.8563 & 1.85 & 0.070 & -7.7961 & 196.2137\\
\bottomrule
\end{tabular}
\end{table}

\begin{table}[H]
\centering
\caption{Estadísticos Globales del Modelo}
\begin{tabular}{lclc}
\toprule
\textbf{Estadístico} & \textbf{Valor} & \textbf{Estadístico} & \textbf{Valor}\\
\midrule
N observaciones & 55      & R-cuadrado      & 0.3698\\
F(1, 53)        & 31.10   & R-cuad. ajustado & 0.3579\\
Prob $>$ F      & 0.0000  & Raíz ECM (RMSE)  & 66.856\\
\bottomrule
\end{tabular}
\end{table}

\subsection{Interpretación}

\begin{cajaconclucion}[title=Interpretación de la Regresión MCO]
\begin{itemize}
  \item \textbf{Intercepto ($\hat{\beta}_0 = 94.21$):} Cuando el gasto total es cero, el gasto estimado en comida sería de 94.21 rupias (sin interpretación económica directa, ya que no se puede tener gasto total igual a cero en la práctica).
  \item \textbf{Pendiente ($\hat{\beta}_1 = 0.4368$):} Por cada rupia adicional de gasto total, el gasto en comida aumenta en promedio \textbf{0.4368 rupias}. Es decir, aproximadamente el 43.7\% del ingreso marginal (representado por el gasto total) se destina a la alimentación.
  \item \textbf{Significancia:} El coeficiente de Gasto\_Total es estadísticamente significativo al 1\% (p-valor = 0.000). El intercepto es marginalmente significativo (p = 0.070).
  \item \textbf{Bondad de ajuste:} El $R^2 = 0.3698$ indica que el gasto total explica el \textbf{36.98\%} de la variación en el gasto en comida, lo cual es razonable para datos de corte transversal.
  \item \textbf{La ecuación estimada es:}
\end{itemize}
\[
\widehat{\text{Gasto\_en\_Comida}} = 94.2088 + 0.4368 \cdot \text{Gasto\_Total}
\]
\end{cajaconclucion}

%===============================================================
\section{Inciso B: Gráfica de Residuos vs.\ Gasto Total}

\subsection{Comandos en Stata}
\begin{cajacomando}[title=Stata -- Generación y Gráfica de Residuos]
\texttt{. predict errores, residuals}\\
\texttt{. gen R = errores\^{}2}\\
\texttt{. twoway (scatter R Gasto\_Total)}
\end{cajacomando}

\subsection{Análisis Visual}

La gráfica de dispersión de los residuos al cuadrado ($\hat{u}^2$) contra el Gasto\_Total (Imagen 1) muestra:

\begin{cajaresultado}[title=Diagnóstico Visual de Heteroscedasticidad]
\begin{itemize}
  \item Para valores bajos de Gasto\_Total (alrededor de 380--500), los residuos al cuadrado son relativamente \textbf{pequeños y concentrados} cerca de cero.
  \item Conforme el Gasto\_Total aumenta (600--800), la dispersión de $\hat{u}^2$ se vuelve \textbf{notablemente mayor}, con valores que alcanzan hasta 30,000.
  \item Este patrón de \textbf{abanico o cono} es la señal clásica de heteroscedasticidad: la varianza de los errores crece con la variable independiente.
  \item \textbf{Conclusión visual:} Se sospecha fuertemente la presencia de heteroscedasticidad. El supuesto $\text{Var}(u_i) = \sigma^2$ (varianza constante) no se satisface.
\end{itemize}
\end{cajaresultado}

%===============================================================
\section{Inciso C: Pruebas Formales de Heteroscedasticidad}

\subsection{Prueba de Breusch-Pagan / Cook-Weisberg (LM)}

\subsubsection{Hipótesis Formales}

\[
H_0: \text{Var}(u_i \mid X_i) = \sigma^2 \quad \text{(varianza constante — homocedasticidad)}
\]
\[
H_a: \text{Var}(u_i \mid X_i) = f(X_i) \quad \text{(varianza depende de } X \text{ — heteroscedasticidad)}
\]

\textbf{Estadístico:} $\text{LM} = n \cdot R^2_{\text{aux}} \sim \chi^2(p)$ bajo $H_0$, donde $p$ es el número de regresores en la regresión auxiliar ($\hat{u}^2$ sobre $\hat{Y}$).

\textbf{Regla de decisión:} Rechazar $H_0$ si LM $> \chi^2_{\alpha}(p)$ o equivalentemente si p-valor $< \alpha$.

\subsubsection{Comando en Stata}
\begin{cajacomando}[title=Stata -- Prueba Breusch-Pagan]
\texttt{. estat hettest}
\end{cajacomando}

\subsubsection{Resultados}
\begin{table}[H]
\centering
\caption{Prueba Breusch-Pagan/Cook-Weisberg}
\begin{tabular}{lcc}
\toprule
\textbf{Hipótesis nula} & \multicolumn{2}{c}{$H_0$: Varianza constante (homocedasticidad)}\\
\midrule
$\chi^2(1)$ & = & 7.18\\
Prob $> \chi^2$ & = & \textbf{0.0074}\\
\bottomrule
\end{tabular}
\end{table}

\subsubsection{Interpretación}
\begin{cajaconclucion}[title=Conclusión -- Breusch-Pagan]
El estadístico $\chi^2(1) = 7.18$ con p-valor $= 0.0074 < 0.05$ (e incluso $< 0.01$). Se \textbf{rechaza la hipótesis nula de varianza constante} al 1\% de significancia. La prueba confirma formalmente la presencia de \textbf{heteroscedasticidad} en el modelo original.

La prueba B-P es de tipo multiplicador de Lagrange (LM): regresa los residuos al cuadrado sobre los valores ajustados y utiliza un test $\chi^2$. Bajo $H_0$: $\text{Var}(u_i \mid X_i) = \sigma^2$. El rechazo implica que la varianza depende sistemáticamente del Gasto Total.
\end{cajaconclucion}

\subsection{Prueba de White}

\subsubsection{Hipótesis Formales}

\[
H_0: \text{Var}(u_i \mid X_i) = \sigma^2 \quad \text{(homocedasticidad)}
\]
\[
H_a: \text{Var}(u_i) \text{ depende de } X_i, X_i^2 \text{ y/o sus productos cruzados (heteroscedasticidad no restringida)}
\]

\textbf{Estadístico:} $W = n \cdot R^2_{\text{aux}} \sim \chi^2(p)$ bajo $H_0$, donde la regresión auxiliar incluye los regresores originales, sus cuadrados y productos cruzados.

\textbf{Regla de decisión:} Rechazar $H_0$ si p-valor $< \alpha$.

\subsubsection{Comando en Stata}
\begin{cajacomando}[title=Stata -- Prueba de White]
\texttt{. estat imtest, white}
\end{cajacomando}

\subsubsection{Resultados}
\begin{table}[H]
\centering
\caption{Prueba de White para Heteroscedasticidad No Restringida}
\begin{tabular}{lcc}
\toprule
\textbf{Fuente} & $\chi^2$ & gl \quad p-valor\\
\midrule
Heteroscedasticidad & 4.14 & 2 \quad 0.1259\\
Asimetría (Skewness) & 7.81 & 1 \quad 0.0052\\
Curtosis (Kurtosis)  & 3.50 & 1 \quad 0.0615\\
\midrule
\textbf{Total (IM-test)} & \textbf{15.45} & 4 \quad \textbf{0.0038}\\
\bottomrule
\end{tabular}
\end{table}

La prueba de White específicamente para heteroscedasticidad arroja:
\[
\chi^2(2) = 4.14, \quad \text{Prob} > \chi^2 = 0.1259
\]

\subsubsection{Interpretación}
\begin{cajaconclucion}[title=Conclusión -- Prueba de White]
\begin{itemize}
  \item La \textbf{prueba de White para heteroscedasticidad} ($\chi^2 = 4.14$, p $= 0.1259 > 0.05$) \textbf{no rechaza} $H_0$ de homocedasticidad. Esto puede parecer contradictorio con B-P, pero la prueba de White con un regresor solo tiene 2 grados de libertad (incluye $X$ y $X^2$) y tiene menos poder en muestras moderadas.
  \item La \textbf{descomposición Cameron-Trivedi} del IM-test muestra que el total es significativo ($\chi^2 = 15.45$, p $= 0.0038$), principalmente impulsado por la \textbf{asimetría} de los residuos ($\chi^2 = 7.81$, p $= 0.0052$), lo cual refuerza la necesidad de transformar el modelo.
  \item En conjunto con B-P, la evidencia apoya la presencia de heteroscedasticidad.
\end{itemize}
\end{cajaconclucion}



%---------------------------------------------------------------
\subsection{Prueba de Park}

\subsubsection{Fundamento Teórico}

Park (1966) propone que la varianza del error sigue la forma $\sigma_i^2 = \sigma^2 X_i^{\beta} e^{v_i}$. Tomando logaritmos:
\[
\ln(\hat{u}_i^2) = \ln(\sigma^2) + \beta \ln(X_i) + v_i = \alpha_0 + \alpha_1 \ln(X_i) + v_i
\]
Si $\alpha_1$ es significativo, existe heteroscedasticidad.

\subsubsection{Hipótesis Formales}

\[
H_0: \alpha_1 = 0 \quad \text{(la varianza no depende de } X \text{ — homocedasticidad)}
\]
\[
H_a: \alpha_1 \neq 0 \quad \text{(la varianza depende de } X^{\alpha_1} \text{ — heteroscedasticidad)}
\]

\textbf{Estadístico:} $t = \hat{\alpha}_1 / \text{SE}(\hat{\alpha}_1) \sim t(n-2)$ bajo $H_0$.

\textbf{Regla de decisión:} Rechazar $H_0$ si $|t| > t_{\alpha/2}(n-2)$ o p-valor $< \alpha$.

\subsubsection{Comandos en Stata}
\begin{cajacomando}[title=Stata -- Prueba de Park]
\texttt{. gen ln\_u2 = ln(errores\^{}2)}\\
\texttt{. gen ln\_x  = ln(Gasto\_Total)}\\
\texttt{. reg ln\_u2 ln\_x}
\end{cajacomando}

\subsubsection{Resultados de Stata}
\begin{table}[H]
\centering
\caption{Prueba de Park: $\ln(\hat{u}^2)$ sobre $\ln(\text{Gasto\_Total})$}
\begin{tabular}{lcccccc}
\toprule
\textbf{Variable} & \textbf{Coef.} & \textbf{Error Est.} & \textbf{t} & \textbf{p-valor} & \textbf{IC 95\% Inf.} & \textbf{IC 95\% Sup.}\\
\midrule
$\ln(\text{Gasto\_Total})$ & 3.7032 & 1.5519 & 2.39 & 0.0206 & 0.5907 & 6.8157\\
\_cons & $-16.8629$ & 10.0014 & $-1.69$ & 0.0977 & $-36.928$ & 3.203\\
\midrule
\multicolumn{3}{l}{$R^2 = 0.0970$, \; $R^2_{\text{adj}} = 0.0800$} & \multicolumn{3}{l}{F(1,53) = 5.69, \; Prob $>$ F = 0.0206}\\
\bottomrule
\end{tabular}
\end{table}

\textbf{Decisión:} $|t| = 2.39 > t_{0.025}(53) \approx 2.006$ y p-valor $= 0.0206 < 0.05$ $\Rightarrow$ \textbf{Se rechaza $H_0$}.

\subsubsection{Interpretación}
\begin{cajaconclucion}[title=Conclusión -- Prueba de Park]
El coeficiente $\hat{\alpha}_1 = 3.7032$ es significativo (p = 0.0206 $< 0.05$). Se rechaza $H_0$. La varianza del error crece aproximadamente como $X^{3.7}$, lo que indica heteroscedasticidad fuerte y creciente. Esta forma funcional ratifica que la transformación logarítmica es una corrección adecuada.
\end{cajaconclucion}

%---------------------------------------------------------------
\subsection{Prueba de Glejser}

\subsubsection{Fundamento Teórico}

Glejser (1969) propone regresar el \textbf{valor absoluto de los residuos} $|\hat{u}_i|$ sobre diferentes formas funcionales de $X_i$. Se prueban tres especificaciones:
\begin{enumerate}
  \item $|\hat{u}_i| = \alpha_0 + \alpha_1 X_i + v_i$ \quad (lineal)
  \item $|\hat{u}_i| = \alpha_0 + \alpha_1 \sqrt{X_i} + v_i$ \quad (raíz cuadrada)
  \item $|\hat{u}_i| = \alpha_0 + \alpha_1 \dfrac{1}{X_i} + v_i$ \quad (inversa)
\end{enumerate}

\subsubsection{Hipótesis Formales}

Para \textbf{cada} forma funcional:
\[
H_0: \alpha_1 = 0 \quad \text{(la varianza no depende de } X \text{ — homocedasticidad)}
\]
\[
H_a: \alpha_1 \neq 0 \quad \text{(la varianza depende de } X \text{ — heteroscedasticidad)}
\]

\textbf{Estadístico:} $t = \hat{\alpha}_1 / \text{SE}(\hat{\alpha}_1) \sim t(n-2)$ bajo $H_0$.\\
\textbf{Regla de decisión:} Rechazar $H_0$ si $|t| > t_{0.025}(53) \approx 2.006$, equivalente a p-valor $< 0.05$.

\subsubsection{Comandos en Stata}
\begin{cajacomando}[title=Stata -- Prueba de Glejser]
\texttt{. gen abs\_u  = abs(errores)}\\
\texttt{. gen sqrt\_x = sqrt(Gasto\_Total)}\\
\texttt{. gen inv\_x  = 1/Gasto\_Total}\\[0.3em]
\texttt{. reg abs\_u Gasto\_Total} \qquad \textit{* Forma 1: lineal}\\
\texttt{. reg abs\_u sqrt\_x} \qquad \quad \textit{* Forma 2: raíz cuadrada}\\
\texttt{. reg abs\_u inv\_x} \qquad \quad \textit{* Forma 3: inversa}
\end{cajacomando}

\subsubsection{Resultados por Forma Funcional}

\medskip
\noindent\textbf{Forma 1 -- Lineal: $|\hat{u}_i| = \alpha_0 + \alpha_1 X_i + v_i$}

\begin{table}[H]
\centering
\caption{Glejser Forma 1: \texttt{reg abs\_u Gasto\_Total}}
\begin{tabular}{lcccccc}
\toprule
\textbf{Variable} & \textbf{Coef.} & \textbf{Error Est.} & \textbf{t} & \textbf{p-valor} & \textbf{IC 95\% Inf.} & \textbf{IC 95\% Sup.}\\
\midrule
Gasto\_Total & 0.1307 & 0.0454 & \textbf{2.88} & \textbf{0.006} & 0.0396 & 0.2218\\
\_cons & $-32.2197$ & 29.4900 & $-1.09$ & 0.280 & $-91.369$ & 26.930\\
\midrule
\multicolumn{3}{l}{$R^2 = 0.1352$, $R^2_{\text{adj}} = 0.1188$} & \multicolumn{3}{l}{F(1,53) = 8.28, \; Prob $>$ F = 0.0058}\\
\bottomrule
\end{tabular}
\end{table}
\textbf{Decisión:} $t = 2.88 > 2.006$ y p = 0.006 $< 0.05$ $\Rightarrow$ \textbf{Se rechaza $H_0$}. $|\hat{u}|$ crece 0.1307 rupias por cada rupia adicional de gasto total.

\medskip
\noindent\textbf{Forma 2 -- Raíz Cuadrada: $|\hat{u}_i| = \alpha_0 + \alpha_1 \sqrt{X_i} + v_i$}

\begin{table}[H]
\centering
\caption{Glejser Forma 2: \texttt{reg abs\_u sqrt\_x}}
\begin{tabular}{lcccccc}
\toprule
\textbf{Variable} & \textbf{Coef.} & \textbf{Error Est.} & \textbf{t} & \textbf{p-valor} & \textbf{IC 95\% Inf.} & \textbf{IC 95\% Sup.}\\
\midrule
sqrt\_x & 6.1994 & 2.2134 & \textbf{2.80} & \textbf{0.007} & 1.7598 & 10.639\\
\_cons & $-104.7165$ & 55.9541 & $-1.87$ & 0.067 & $-216.946$ & 7.513\\
\midrule
\multicolumn{3}{l}{$R^2 = 0.1289$, $R^2_{\text{adj}} = 0.1125$} & \multicolumn{3}{l}{F(1,53) = 7.84, \; Prob $>$ F = 0.0071}\\
\bottomrule
\end{tabular}
\end{table}
\textbf{Decisión:} $t = 2.80 > 2.006$ y p = 0.007 $< 0.05$ $\Rightarrow$ \textbf{Se rechaza $H_0$}. Sugiere $\sigma_i^2 \propto X_i$ (varianza proporcional al gasto total).

\medskip
\noindent\textbf{Forma 3 -- Inversa: $|\hat{u}_i| = \alpha_0 + \alpha_1 (1/X_i) + v_i$}

\begin{table}[H]
\centering
\caption{Glejser Forma 3: \texttt{reg abs\_u inv\_x}}
\begin{tabular}{lcccccc}
\toprule
\textbf{Variable} & \textbf{Coef.} & \textbf{Error Est.} & \textbf{t} & \textbf{p-valor} & \textbf{IC 95\% Inf.} & \textbf{IC 95\% Sup.}\\
\midrule
inv\_x & $-37567.74$ & 15057.74 & $\mathbf{-2.49}$ & \textbf{0.016} & $-67769.75$ & $-7365.72$\\
\_cons & 112.4307 & 25.0693 & 4.48 & 0.000 & 62.148 & 162.713\\
\midrule
\multicolumn{3}{l}{$R^2 = 0.1051$, $R^2_{\text{adj}} = 0.0882$} & \multicolumn{3}{l}{F(1,53) = 6.22, \; Prob $>$ F = 0.0157}\\
\bottomrule
\end{tabular}
\end{table}
\textbf{Decisión:} $|t| = 2.49 > 2.006$ y p = 0.016 $< 0.05$ $\Rightarrow$ \textbf{Se rechaza $H_0$}. Signo negativo: $|\hat{u}|$ decrece con $1/X$, es decir, crece con $X$. Consistente con las formas anteriores.

\subsubsection{Gráficos de las Formas de Glejser y Prueba de Park}

\begin{figure}[H]
  \centering
  \includegraphics[width=\textwidth]{grafico_glejser_park.png}
  \caption{Prueba de Glejser (3 formas) y Park: $|\hat{u}_i|$ y $\ln(\hat{u}_i^2)$ contra transformaciones de Gasto\_Total. En los cuatro paneles el coeficiente $\alpha_1$ es significativo (p~$<$~0.05), rechazando $H_0$ de homocedasticidad.}
\end{figure}

\begin{figure}[H]
  \centering
  \includegraphics[width=0.95\textwidth]{grafico_residuos_glejser.png}
  \caption{Izquierda: $\hat{u}_i^2$ vs.\ Gasto Total (patrón de abanico). Derecha: $|\hat{u}_i|$ con las tres líneas ajustadas de Glejser --- todas crecientes, confirmando heteroscedasticidad.}
\end{figure}

\subsubsection{Resumen Glejser -- Las Tres Formas}
\begin{table}[H]
\centering
\caption{Síntesis Prueba de Glejser}
\begin{tabular}{lccccl}
\toprule
\textbf{Forma} & $\hat{\alpha}_1$ & $t$ & \textbf{p-valor} & $R^2$ & \textbf{Decisión}\\
\midrule
1. Lineal   & $0.1307$ & $2.88$ & $0.006$ & $0.135$ & Rechazar $H_0$ \\
2. Raíz     & $6.1994$ & $2.80$ & $0.007$ & $0.129$ & Rechazar $H_0$ \\
3. Inversa  & $-37567.74$ & $-2.49$ & $0.016$ & $0.105$ & Rechazar $H_0$ \\
\bottomrule
\end{tabular}
\end{table}

\begin{cajaconclucion}[title=Conclusión -- Prueba de Glejser]
\begin{itemize}
  \item Las \textbf{tres formas funcionales} rechazan $H_0$ (p $<$ 0.05). La heteroscedasticidad es robusta a la especificación elegida.
  \item La \textbf{Forma 1 lineal} tiene el mejor ajuste ($R^2 = 0.135$) y la mayor significancia (p = 0.006), siendo la especificación preferida.
  \item La \textbf{Forma 2 de raíz cuadrada} (p = 0.007) sugiere que la desviación estándar de los errores crece con $\sqrt{X}$, es decir $\sigma_i \propto \sqrt{X_i}$, lo que implica $\sigma_i^2 \propto X_i$. Esta es la forma de heteroscedasticidad más común en datos de gasto familiar.
  \item La \textbf{Forma 3 inversa} (p = 0.016) confirma que $|\hat{u}|$ crece con $X$ (el coeficiente negativo de $1/X$ es equivalente a una relación positiva con $X$).
\end{itemize}
\end{cajaconclucion}

%---------------------------------------------------------------
\subsection{Prueba de Goldfeld-Quandt}

\subsubsection{Fundamento Teórico}

La prueba de Goldfeld-Quandt (1965) supone que la varianza de los errores es mayor para observaciones con $X$ alto que para observaciones con $X$ bajo. El procedimiento es:
\begin{enumerate}
  \item Ordenar las observaciones por la variable independiente $X$ (Gasto\_Total). \textbf{Los datos de la Tabla 2.8 ya están ordenados de forma ascendente por Gasto\_Total.}
  \item Eliminar $c$ observaciones del centro para separar claramente los dos grupos.
  \item Estimar MCO por separado en cada subgrupo y obtener su Suma de Cuadrados de los Residuos (SCR).
  \item Calcular el estadístico $F = \text{SCR}_2 / \text{SCR}_1$ (el grupo con $X$ alto sobre el grupo con $X$ bajo).
  \item Bajo $H_0$ de homocedasticidad, $F \sim F\left(\frac{n-c}{2}-k,\; \frac{n-c}{2}-k\right)$.
\end{enumerate}

\subsubsection{Hipótesis Formales}

\[
H_0: \sigma_1^2 = \sigma_2^2 \quad \text{(homocedasticidad: igual varianza en ambos grupos)}
\]
\[
H_a: \sigma_2^2 > \sigma_1^2 \quad \text{(la varianza del grupo de gasto alto es mayor — heteroscedasticidad)}
\]

\textbf{Estadístico:} $\lambda = \dfrac{\text{SCR}_2/(m-k)}{\text{SCR}_1/(m-k)} = \dfrac{\text{SCR}_2}{\text{SCR}_1} \sim F(m-k,\; m-k)$ bajo $H_0$.\\[0.3em]
\textbf{Regla de decisión:} Rechazar $H_0$ si $\lambda > F_{\alpha}(m-k,\; m-k)$ (prueba unilateral derecha).

\subsubsection{Definición de los Grupos ($c = 5$, $n = 55$)}

Con $n = 55$ observaciones y eliminando $c = 5$ observaciones centrales, cada grupo queda con:
\[
m = \frac{n - c}{2} = \frac{55 - 5}{2} = \frac{50}{2} = 25 \text{ observaciones}
\]

\begin{table}[H]
\centering
\caption{Partición de Observaciones -- Prueba Goldfeld-Quandt}
\begin{tabular}{lccc}
\toprule
\textbf{Grupo} & \textbf{Observaciones} & \textbf{Gasto Total (rango)} & \textbf{n por grupo}\\
\midrule
Grupo 1 (bajo)    & obs 1 -- 25  & 382 -- 635 rupias & 25\\
\textit{Eliminadas} & \textit{obs 26 -- 30} & \textit{640 -- 662 rupias} & \textit{5}\\
Grupo 2 (alto)    & obs 31 -- 55 & 663 -- 801 rupias & 25\\
\bottomrule
\end{tabular}
\end{table}

\subsubsection{Comandos en Stata}

\begin{cajacomando}[title=Stata -- Prueba de Goldfeld-Quandt]
\texttt{* Grupo 1: observaciones 1 a 25 (Gasto\_Total bajo)}\\
\texttt{. reg Gasto\_en\_Comida Gasto\_Total if \_n <= 25}\\[0.4em]
\texttt{* Grupo 2: observaciones 31 a 55 (Gasto\_Total alto)}\\
\texttt{. reg Gasto\_en\_Comida Gasto\_Total if \_n >= 31}\\[0.4em]
\texttt{* Alternativamente, usando la variable Gasto\_Total directamente:}\\
\texttt{. reg Gasto\_en\_Comida Gasto\_Total if Gasto\_Total <= 635}\\
\texttt{. reg Gasto\_en\_Comida Gasto\_Total if Gasto\_Total >= 663}
\end{cajacomando}

\subsubsection{Resultados de Stata y Cálculo del Estadístico}

Con $k = 2$ parámetros y $m = 25$ observaciones por grupo, los grados de libertad son $m - k = 23$ por grupo.

\begin{table}[H]
\centering
\caption{Resultados de las Regresiones por Grupo -- Goldfeld-Quandt}
\begin{tabular}{lcc}
\toprule
\textbf{Estadístico} & \textbf{Grupo 1 (obs 1--25, bajo)} & \textbf{Grupo 2 (obs 31--55, alto)}\\
\midrule
$n$ por grupo      & 25        & 25\\
Gasto\_Total rango & 382--635  & 663--801\\
$\hat{\beta}_0$    & 65.7368   & 420.6548\\
$\hat{\beta}_1$    & 0.4921    & 0.0014\\
$R^2$              & 0.4917    & 0.0000\\
F del modelo       & 22.25 (p=0.0001) & 0.00 (p=0.9973)\\
Raíz ECM           & 43.389    & 83.269\\
\textbf{SCR (Residual SS)} & \textbf{43{,}300.21} & \textbf{159{,}475.36}\\
gl residual        & 23        & 23\\
\bottomrule
\end{tabular}
\end{table}

El estadístico de la prueba es:
\[
\lambda = \frac{\text{SCR}_2 / (m-k)}{\text{SCR}_1 / (m-k)} = \frac{\text{SCR}_2}{\text{SCR}_1} = \frac{159{,}475.36}{43{,}300.21} = \mathbf{3.683}
\]

\subsubsection{Hipótesis y Regla de Decisión}

\[
H_0: \sigma_1^2 = \sigma_2^2 \quad \text{(homocedasticidad)}
\]
\[
H_a: \sigma_2^2 > \sigma_1^2 \quad \text{(el grupo de alto gasto tiene mayor varianza)}
\]

\textbf{Distribución bajo $H_0$:} $\lambda \sim F(23,\; 23)$

\textbf{Valor crítico:} $F_{0.05}(23, 23) = 2.014$ \quad (tabla F, unilateral derecha, 5\%)

\textbf{p-valor:} $P[F(23,23) > 3.683] = 0.0014$

\textbf{Regla de decisión:}
\[
\lambda = 3.683 > F_{0.05}(23,23) = 2.014 \quad \Rightarrow \quad \textbf{Se rechaza } H_0
\]

\subsubsection{Interpretación}

\begin{cajaconclucion}[title=Conclusión -- Prueba de Goldfeld-Quandt]
\begin{itemize}
  \item $\lambda = 3.683$ con p-valor $= 0.0014 < 0.01$. Se rechaza $H_0$ al \textbf{1\% de significancia}: la varianza del grupo de gasto alto es estadísticamente mayor que la del grupo de gasto bajo.

  \item \textbf{Nota sobre el Grupo 2:} El $R^2 = 0.0000$ y el coeficiente de Gasto\_Total prácticamente nulo (0.0014) en el grupo de gasto alto no es un error de datos --- refleja que en el rango 663--801 rupias, el gasto en comida presenta una dispersión tan elevada y asistemática que no hay relación lineal detectable. Esto es precisamente la manifestación de la heteroscedasticidad: la varianza es tan grande en ese tramo que ``ahoga'' la señal de la pendiente. La raíz ECM del Grupo 2 (83.27) prácticamente duplica la del Grupo 1 (43.39), y SCR$_2 \approx 3.7 \times$ SCR$_1$.

  \item \textbf{Consistencia con las demás pruebas:} Breusch-Pagan (p = 0.0074), regresión auxiliar Park (p = 0.011) y ahora Goldfeld-Quandt ($\lambda = 3.683$, p = 0.0014) apuntan en la misma dirección: existe heteroscedasticidad significativa que crece con el nivel de gasto total.

  \item \textbf{Ventaja de GQ:} intuitiva, basada en distribución F exacta bajo normalidad. \textbf{Desventaja:} sensible a la elección de $c$ y pierde $c$ observaciones. Con $c = 5$ y $n = 55$ la separación es suficiente sin sacrificar demasiada información.
\end{itemize}
\end{cajaconclucion}

%===============================================================
\section{Inciso D: Corrección mediante Transformación Logarítmica}

\subsection{Justificación Teórica}

Cuando la varianza de los errores crece con $X$ (posiblemente de forma proporcional), una transformación logarítmica de ambas variables frecuentemente estabiliza la varianza. El modelo log-log es:
\[
\ln(\text{Gasto\_en\_Comida}_i) = \beta_0 + \beta_1 \ln(\text{Gasto\_Total}_i) + u_i
\]
En este modelo, $\beta_1$ se interpreta como la \textbf{elasticidad}: el cambio porcentual en el gasto en comida ante un cambio del 1\% en el gasto total.

\subsection{Comandos en Stata}
\begin{cajacomando}[title=Stata -- Transformación Logarítmica]
\texttt{. gen logaritmo\_y = ln(Gasto\_en\_Comida)}\\
\texttt{. gen logaritmo\_x = ln(Gasto\_Total)}\\
\texttt{. reg logaritmo\_y logaritmo\_x}\\
\texttt{. estat hettest}
\end{cajacomando}

\subsection{Resultados}

\begin{table}[H]
\centering
\caption{Regresión Log-Log: logaritmo\_y sobre logaritmo\_x}
\begin{tabular}{lcccccc}
\toprule
\textbf{Variable} & \textbf{Coef.} & \textbf{Err. Est.} & \textbf{t} & \textbf{P$>|t|$} & \textbf{IC 95\% Inf.} & \textbf{IC 95\% Sup.}\\
\midrule
logaritmo\_x & 0.7363 & 0.1207 & 6.10 & 0.000 & 0.4942 & 0.9784\\
\_cons       & 1.1543 & 0.7780 & 1.48 & 0.144 & -0.4061 & 2.7147\\
\bottomrule
\end{tabular}
\end{table}

\begin{table}[H]
\centering
\caption{Estadísticos Globales -- Modelo Log-Log}
\begin{tabular}{lclc}
\toprule
N observaciones & 55   & R-cuadrado       & 0.4125\\
F(1, 53)        & 37.21 & R-cuad. ajustado & 0.4014\\
Prob $>$ F      & 0.0000 & Raíz ECM        & 0.17662\\
\bottomrule
\end{tabular}
\end{table}

\subsection{Verificación de Heteroscedasticidad en el Modelo Log-Log}

\begin{table}[H]
\centering
\caption{Breusch-Pagan en el Modelo Log-Log}
\begin{tabular}{lcc}
\toprule
$\chi^2(1)$ & = & 1.60\\
Prob $> \chi^2$ & = & \textbf{0.2052}\\
\bottomrule
\end{tabular}
\end{table}

\subsection{Interpretación}

\begin{cajaconclucion}[title=Interpretación del Modelo Log-Log]
\begin{enumerate}
  \item \textbf{Elasticidad ($\hat{\beta}_1 = 0.7363$):} Por cada 1\% de aumento en el gasto total, el gasto en comida aumenta en promedio \textbf{0.7363\%}. El bien ``alimentos'' se comporta como un \textbf{bien normal no de lujo} (elasticidad $< 1$), lo cual es consistente con la Ley de Engel: la proporción del ingreso destinada a alimentos disminuye conforme el ingreso aumenta.
  
  \item \textbf{Mejora en el ajuste:} El $R^2$ mejora de 0.3698 (modelo lineal) a \textbf{0.4125} (modelo log-log), indicando que la especificación logarítmica captura mejor la relación.
  
  \item \textbf{Corrección de la heteroscedasticidad:} El test de Breusch-Pagan en el modelo log-log da $\chi^2(1) = 1.60$ con p-valor $= 0.2052 > 0.05$. \textbf{No se rechaza $H_0$} de varianza constante. La transformación logarítmica \textbf{corrigió exitosamente} el problema de heteroscedasticidad.
  
  \item \textbf{La ecuación estimada (en logaritmos) es:}
  \[
  \widehat{\ln(\text{Gasto\_en\_Comida})} = 1.1543 + 0.7363 \cdot \ln(\text{Gasto\_Total})
  \]
  \[
  \text{o equivalentemente:} \quad \widehat{\text{Gasto\_en\_Comida}} = e^{1.1543} \cdot \text{Gasto\_Total}^{0.7363}
  \]
\end{enumerate}
\end{cajaconclucion}

%===============================================================
\section{Inciso E: Errores Estándar Robustos de White}

\subsection{Fundamento Teórico}

Los \textbf{errores estándar robustos} (consistentes con la heteroscedasticidad) propuestos por \textbf{White (1980)} permiten realizar inferencia válida incluso cuando existe heteroscedasticidad, sin necesidad de corregir el modelo. La varianza corregida es:
\[
\widehat{\text{Var}}_{HC}(\hat{\boldsymbol{\beta}}) = (X'X)^{-1} \left(\sum_{i=1}^{n} \hat{u}_i^2 x_i x_i'\right) (X'X)^{-1}
\]

\subsection{Comando en Stata}
\begin{cajacomando}[title=Stata -- Regresión con Errores Robustos de White]
\texttt{. reg Gasto\_en\_Comida Gasto\_Total, robust}
\end{cajacomando}

\subsection{Resultados}

\begin{table}[H]
\centering
\caption{Regresión con Errores Robustos (White): Gasto\_en\_Comida sobre Gasto\_Total}
\begin{tabular}{lcccccc}
\toprule
\textbf{Variable} & \textbf{Coef.} & \textbf{Err. Robusto} & \textbf{t} & \textbf{P$>|t|$} & \textbf{IC 95\% Inf.} & \textbf{IC 95\% Sup.}\\
\midrule
Gasto\_Total & 0.4368 & 0.0743 & 5.88 & 0.000 & 0.2879 & 0.5857\\
\_cons       & 94.2088 & 43.2631 & 2.18 & 0.034 & 7.434 & 180.984\\
\midrule
\multicolumn{3}{l}{F(1,53) = 34.60} & \multicolumn{4}{l}{Prob $>$ F = 0.0000, $R^2 = 0.3698$}\\
\bottomrule
\end{tabular}
\end{table}

\subsection{Comparación de Errores Estándar}

\begin{table}[H]
\centering
\caption{Comparación: Errores MCO vs.\ Errores Robustos de White}
\begin{tabular}{lccc}
\toprule
\textbf{Variable} & \textbf{Coef.} & \textbf{Error Est. MCO} & \textbf{Error Est. Robusto}\\
\midrule
Gasto\_Total & 0.4368 & 0.0783 & \textbf{0.0743}\\
\_cons       & 94.2088 & 50.8563 & \textbf{43.2631}\\
\bottomrule
\end{tabular}
\end{table}

\subsection{Interpretación y Decisión}

\begin{cajaconclucion}[title=Conclusión -- Errores Robustos y Decisión Final]
\begin{enumerate}
  \item \textbf{Diferencias entre errores:} Los errores estándar robustos son \textit{ligeramente menores} que los MCO ordinarios (0.0743 vs.\ 0.0783 para la pendiente; 43.26 vs.\ 50.86 para el intercepto). Esto indica que los errores MCO \textit{sobreestiman} modestamente la incertidumbre en este caso particular.
  
  \item \textbf{Cambio en las conclusiones inferenciales:}
  \begin{itemize}
    \item Gasto\_Total: sigue siendo significativo al 1\% con ambos métodos (p = 0.000 en ambos).
    \item Intercepto: cambia su significancia --- con MCO es marginalmente significativo (p = 0.070), mientras que con errores robustos \textbf{pasa a ser significativo al 5\%} (p = 0.034).
    \item La magnitud de los coeficientes \textbf{no cambia} (como es esperado: la robustez solo afecta los errores estándar, no los coeficientes).
  \end{itemize}
  
  \item \textbf{¿Vale la pena corregir el modelo?}
  
  En este caso se \textbf{recomienda sí corregir}, pero la elección entre métodos depende del objetivo:
  \begin{itemize}
    \item \textbf{Si se desea interpretación directa en niveles}: usar errores robustos de White es suficiente. Los coeficientes son los mismos MCO, pero la inferencia es válida bajo heteroscedasticidad. Esto es la solución más sencilla.
    \item \textbf{Si se desea mayor eficiencia y mejor ajuste}: el modelo log-log es preferible, pues corrige la heteroscedasticidad en la especificación del modelo, mejora el $R^2$ (0.41 vs.\ 0.37) y ofrece una interpretación económicamente significativa (elasticidad).
    \item Dado que la heterocedasticidad fue confirmada significativamente por B-P (p = 0.0074) y la transformación logarítmica la eliminó, se recomienda \textbf{el modelo log-log como la especificación final preferida}.
  \end{itemize}
\end{enumerate}
\end{cajaconclucion}

%===============================================================
\section{Resumen Ejecutivo}

\begin{table}[H]
\centering
\caption{Resumen de Resultados por Inciso}
\begin{tabular}{p{1cm}p{5.5cm}p{7cm}}
\toprule
\textbf{Inc.} & \textbf{Procedimiento} & \textbf{Conclusión Principal}\\
\midrule
A & Regresión MCO lineal & $\hat{Y} = 94.21 + 0.437X$; $R^2 = 0.37$; Gasto Total significativo (p$<$0.001)\\
\addlinespace
B & Gráfica $\hat{u}^2$ vs.\ $X$ & Patrón de abanico sugiere heterocedasticidad creciente con $X$\\
\addlinespace
C & B-P: $\chi^2=7.18$ (p=0.007); White: p=0.126; Park: $\hat{\alpha}_1=3.70$ (p=0.021); Glejser: 3 formas p$<$0.016; GQ: $\lambda=3.683>F_{0.05}(23,23)=2.014$ & Todas las pruebas (excepto White aislado) confirman heteroscedasticidad creciente con $X$\\
\addlinespace
D & Modelo log-log; B-P post: $\chi^2=1.60$ (p=0.205) & Transformación log corrige la heterocedasticidad; elasticidad = 0.74; $R^2$ mejora a 0.41\\
\addlinespace
E & Errores robustos de White & Diferencias pequeñas vs.\ MCO; intercepto cambia significancia; corrección recomendable; modelo log-log es la opción preferida\\
\bottomrule
\end{tabular}
\end{table}

\vspace{2em}
\begin{cajaresultado}[title=Nota sobre los Comandos de Stata]
\textbf{Corrección requerida:} En la imagen se observa el intento de usar \texttt{estat imtest, whiitw} (con doble ``i''), que Stata rechazó con el mensaje \textit{``option whiitw not allowed''}. El comando correcto es:
\[
\texttt{. estat imtest, white}
\]
(con una sola ``i'' en ``white''). Esta corrección está reflejada en la Imagen 4 donde el comando se ejecuta correctamente.
\end{cajaresultado}

\vspace{1em}
\noindent\textbf{Fuente de los datos:} Chandan Mukherjee, Howard White y Marc Wuyts, \textit{Econometrics and Data Analysis for Developing Countries}, Routledge, Nueva York, 1998, p.~457.

\end{document}
